Fix an arm \(a\in\{0,1\}\) and condition on \(\mathcal T\). By \(\texttt{lem:orbit-score-factorization}\), for every Borel set \(C\subseteq[0,\infty)\),
\[
\Pr(S_{i,a}\in C\mid\mathcal T)=\sum_{O\subseteq\mathcal U_a}q_{a,O}F_{a,O}(C),\qquad
\Pr(S_{0,a}\in C\mid\mathcal T)=\sum_{O\subseteq\mathcal U_a}p_{a,O}F_{a,O}(C).\tag{1}
\]
Only the arm-\(a\) orbits have nonzero mass in these sums. If joint orbit balance holds, then \(q_{a,O}=p_{a,O}\) for every arm and orbit by \(\texttt{def:joint-orbit-balance}\). Equation \(1\) therefore proves equality in law of \(S_{i,a}\) and \(S_{0,a}\) for every \(P\in\mathcal P_{\mathcal G,K}\) and every fixed measurable predictor \(\widehat m_a\). The factorization lemma also states that the \(M\) calibration orbit--score pairs are iid and independent of the target pair. Their score marginals now coincide, so \(S_{1,a},\ldots,S_{M,a},S_{0,a}\) are iid conditional on \(\mathcal T\). This reasoning is armwise and does not assert independence or exchangeability between the two arms.

Conversely, suppose joint orbit balance fails. Then there are an arm \(a\) and an orbit \(O_\star\subseteq\mathcal U_a\) such that
\[
q_{a,O_\star}\ne p_{a,O_\star}.\tag{2}
\]
Use \(\texttt{def:orbit-separating-law}\) and the fixed predictors \(\widehat m_0=\widehat m_1=0\). Its arm-\(a\) potential outcome is \(Y_v(a)=K\mathbf 1\{(v,a)\in O_\star\}\), and its other-arm potential outcome is zero. Hence all potential outcomes lie in \([-K,K]\). Let \(L^\star\) denote the prototype and let \(G_0\) be uniform on the finite group \(\mathcal G\). For every \(h\in\mathcal G\), the random elements \(hG_0L^\star\) and \(G_0L^\star\) have the same law because left multiplication permutes the uniform group law. The symmetrized latent law is therefore \(\mathcal G\)-invariant and, by boundedness, belongs to \(\mathcal P_{\mathcal G,K}\). Independent copies give the replicated arrays required by \(\texttt{ass:independent-replicated-arrays}\), and assignments are subsequently drawn from the declared known randomization law.

For every assignment \(z\) with \(\nu(z)>0\), \(\texttt{ass:observable-selector-kernel}\) makes the selector a normalized kernel supported on members satisfying \(z_v=a\). Therefore its selected outcome is observed in arm \(a\), and
\[
S_{i,a}=K\mathbf 1\{U_{i,a}\in O_\star\},\qquad
S_{0,a}=K\mathbf 1\{T^a\in O_\star\}.\tag{3}
\]
Assignments outside \(\operatorname{supp}(\nu)\) have zero weight in the integrated selector law, so their arbitrary kernel extensions cannot alter either identity or any orbit mass. Since \(K>0\), equations \(2\)--\(3\) yield
\[
\Pr(S_{i,a}=K\mid\mathcal T)=q_{a,O_\star}\ne p_{a,O_\star}=\Pr(S_{0,a}=K\mid\mathcal T).
\]
Thus uniform equality of the armwise score laws forces every orbit mass to agree, which is joint orbit balance. The argument imposes no additional relabelling restriction on the selector kernels.

\paragraph{Comparator scope.} This proof uses finite-group invariance to identify common within-orbit score laws, but its conclusion is an equality-of-calibration-and-deployment-law characterization for iid replicated latent arrays with assignment-based calibration selection. The finite permutation and group-testing results of \cite{HemerikGoeman2018,RamdasBarberCandesTibshirani2023} and predictive inference under group symmetry in \cite{DobribanYu2025} provide symmetry comparators; the additional statement here is the universal orbit-mass if-and-only-if and its explicit bounded orbit-separating necessity witness in the present matched-randomization experiment. The causal and counterfactual conformal setting of \cite{LeiCandes2021} and the cluster-randomized conformal setting of \cite{WangLiYu2025} are target and design comparators. These comparisons make no priority or absence claim.